\section*{Capítulo \ref{ch:b1:isostasy-gravity} --- Isostasia y anomalías de la gravedad}

\begin{solution}{exo:b1:isostasy-gravity:1}
(a) $\qty{-150}{\milli\gal} = -150\times\qty{e-5}{m.s^{-2}} = \qty{-1.5e-3}{m.s^{-2}}$,
y $\num{1.5e-3}/9.81 = \num{1.53e-4}$ de $g$. (b) $\qty{1}{\milli\gal} =
\qty{e-5}{m.s^{-2}}$, \num{1.02e-6} de $g$; $\qty{1}{\micro\gal} =
\qty{e-8}{m.s^{-2}}$, \num{1.02e-9} de $g$: un gravímetro que resuelve un microgal mide la gravedad con una precisión de una parte en mil millones.
\end{solution}

\begin{solution}{exo:b1:isostasy-gravity:2}
(a) En el ecuador se anulan ambos senos: $\gamma = \qty{9.780327}{m.s^{-2}}$. A
\ang{45}, $\sin^2\varphi = 0.5$ y $\sin^2 2\varphi = 1$:
$\gamma = 9.780327\times(1 + 0.0026512 - 0.0000058) = \qty{9.806200}{m.s^{-2}}$.
En los polos, $\sin^2\varphi = 1$ y $\sin 2\varphi = 0$:
$\gamma = 9.780327\times1.0053024 = \qty{9.832186}{m.s^{-2}}$.
(b) $9.832186 - 9.780327 = \qty{0.051859}{m.s^{-2}}$, es decir
\qty{5186}{\milli\gal}. (c) $3392/5186 = 0.65$: el efecto centrífugo explica
casi dos tercios de la diferencia; el achatamiento, que acerca los polos al
centro, el resto.
\end{solution}

\begin{solution}{exo:b1:isostasy-gravity:3}
$r = h\rho_c/(\rho_m - \rho_c) = 3\times2800/500 = \qty{16.8}{km}$. La \omterm{def:b1:layered-earth:layers}{corteza} es
más gruesa en $h + r = \qty{19.8}{km}$ (es decir $h\rho_m/(\rho_m - \rho_c) =
3\times3300/500$).
\end{solution}

\begin{solution}{exo:b1:isostasy-gravity:4}
(a) $1/0.3086 = \qty{3.24}{m}$. (b) $0.01/0.3086 = \qty{0.032}{m}$: las alturas
deben conocerse a unos \qty{3}{cm}, por lo que un levantamiento registra la altura de cada estación con posicionamiento por satélite o nivelación.
\end{solution}

\begin{solution}{exo:b1:isostasy-gravity:5}
(a) $\rho = 2700\times100/(100 + h)$: \qty{2673}{kg.m^{-3}} para \qty{1}{km},
\qty{2621}{kg.m^{-3}} para \qty{3}{km}, \qty{2571}{kg.m^{-3}} para \qty{5}{km}.
(b) $\rho = (2700\times100 - 1030\times4)/96 = \num{265880}/96 =
\qty{2770}{kg.m^{-3}}$: la columna bajo el mar es más densa. (c) El modelo de Pratt
exige cambios laterales de densidad a profundidad constante. Bajo el fondo oceánico
tales cambios existen: la \omterm{def:b1:layered-earth:lithosphere}{litosfera} joven y caliente junto a una dorsal es más ligera que la
\omterm{def:b1:layered-earth:lithosphere}{litosfera} vieja y fría alejada de ella (\cref{ch:b1:heat-budget}). Bajo los continentes,
en cambio, el espesor de la \omterm{def:b1:layered-earth:layers}{corteza} se ve que cambia con el relieve.
\end{solution}

\begin{solution}{exo:b1:isostasy-gravity:6}
(a) Estación 1: $g - \gamma = 9.807751 - 9.808911 = \qty{-0.001160}{m.s^{-2}} =
\qty{-116.0}{\milli\gal}$; \omterm{def:b1:isostasy-gravity:anomalies}{corrección de aire libre} $0.3086\times350 =
\qty{108.0}{\milli\gal}$, así que $\Delta g_{\mathrm{FA}} = \qty{-8.0}{\milli\gal}$;
\omterm{def:b1:isostasy-gravity:anomalies}{corrección de Bouguer} $0.112\times350 = \qty{39.2}{\milli\gal}$, así que
$\Delta g_{\mathrm B} = \qty{-47.2}{\milli\gal}$.
Estación 2: $g - \gamma = \qty{-535.9}{\milli\gal}$; \omterm{def:b1:isostasy-gravity:anomalies}{corrección de aire libre}
$0.3086\times1850 = \qty{570.9}{\milli\gal}$, así que $\Delta g_{\mathrm{FA}} =
\qty{+35.0}{\milli\gal}$; \omterm{def:b1:isostasy-gravity:anomalies}{corrección de Bouguer} $0.112\times1850 =
\qty{207.2}{\milli\gal}$, así que $\Delta g_{\mathrm B} = \qty{-172.2}{\milli\gal}$.
(b) En la estación 2 la \omterm{def:b1:isostasy-gravity:anomalies}{anomalía de Bouguer} es fuertemente negativa: bajo las montañas
hay un déficit de masa, una raíz de \omterm{def:b1:layered-earth:layers}{corteza} ligera. La \omterm{def:b1:isostasy-gravity:anomalies}{anomalía de aire libre}, pequeña
comparada con la \omterm{def:b1:isostasy-gravity:anomalies}{corrección de Bouguer}, muestra que el déficit casi equilibra las montañas:
están en gran parte compensadas (en torno a $172/207$, o \qty{83}{\percent}, si el modelo de
placa vale para toda la cordillera).
\end{solution}

\begin{solution}{exo:b1:isostasy-gravity:7}
(a) Como en la prueba de la proposición, una capa delgada situada a la distancia $z$
contribuye $G\rho\,2\pi z\,h\int_0^a s\,(s^2 + z^2)^{-3/2}\,\mathrm ds$, ahora
con la integral limitada a $s = a$:
$\int_0^a s\,(s^2 + z^2)^{-3/2}\,\mathrm ds = 1/z - 1/\sqrt{z^2 + a^2}$. Así
$g = 2\pi G\rho h\,\bigl(1 - z/\sqrt{z^2 + a^2}\bigr)$, que tiende al valor de la placa
cuando $a \to \infty$. (b) $z/\sqrt{z^2 + a^2} = 0.1$ da
$\sqrt{z^2 + a^2} = 10z$, $a = \sqrt{99}\,z = \qty{9.95}{km}$ para
$z = \qty{1}{km}$. (c) La roca dentro de unas diez veces su profundidad bajo la
estación aporta el \qty{90}{\percent} de la atracción: la corrección del terreno
necesita el relieve en un radio de varios kilómetros a decenas de kilómetros alrededor
de cada estación, no solo bajo ella.
\end{solution}

\begin{solution}{exo:b1:isostasy-gravity:8}
(a) $r' = d(\rho_c - \rho_w)/(\rho_m - \rho_c) = 4\times1770/500 =
\qty{14.16}{km}$, y la \omterm{def:b1:layered-earth:layers}{corteza} es $35 - 4 - 14.16 = \qty{16.84}{km}$ de espesor.
(b) El modelo de Airy con una sola densidad cortical predice una \omterm{def:b1:layered-earth:layers}{corteza} oceánica más de
tres veces demasiado gruesa. La \omterm{def:b1:layered-earth:layers}{corteza} oceánica no es \omterm{def:b1:layered-earth:layers}{corteza} continental adelgazada: es
una roca distinta y más densa, y los dos niveles de la Tierra son dos tipos de \omterm{def:b1:layered-earth:layers}{corteza},
como sugiere la \omterm{def:b1:isostasy-gravity:hypsometry}{curva hipsométrica}.
\end{solution}

\begin{solution}{exo:b1:isostasy-gravity:9}
(a) Un segundo de arco es $\pi/(180\times3600) = 1/\num{206265}$ de radián:
\qty{5.236}{''} $= \qty{2.538e-5}{rad}$ y \qty{15.885}{''} $=
\qty{7.701e-5}{rad}$. (b) $a_h = g\theta$: $9.79\times\num{2.538e-5} =
\qty{2.485e-4}{m.s^{-2}} = \qty{24.9}{\milli\gal}$ observado, y
$9.79\times\num{7.701e-5} = \qty{75.4}{\milli\gal}$ calculado a partir de las
montañas visibles. (c) $1 - 5.236/15.885 = 0.670$: la compensación canceló cerca de
dos tercios de la atracción de las montañas visibles.
\end{solution}

\begin{solution}{exo:b1:isostasy-gravity:10}
(a) Fuera de ella, una esfera atrae como una masa puntual en su centro, con
$GM/(x^2 + z^2)$ a lo largo de la recta hasta el centro; la componente vertical es eso
por el coseno $z/\sqrt{x^2 + z^2}$, que da $GMz/(x^2 + z^2)^{3/2}$.
(b) $M = \tfrac43\pi(\qty{2000}{m})^3\times(\qty{-200}{kg.m^{-3}}) =
\qty{-6.70e12}{kg}$, y sobre el centro $GM/z^2 = \num{6.674e-11}\times
(\num{-6.70e12})/(4000)^2 = \qty{-2.80e-5}{m.s^{-2}} = \qty{-2.80}{\milli\gal}$.
(c) $(1 + x^2/z^2)^{-3/2} = 1/2$ da $x = z\sqrt{2^{2/3} - 1} = 0.766\,z =
\qty{3.07}{km}$. (d) La semianchura a media altura del perfil da la profundidad:
$z = 1.305\,x_{1/2}$. Un cuerpo más profundo da una anomalía más ancha y más suave.
\end{solution}

\begin{solution}{exo:b1:isostasy-gravity:11}
(a) $V = \tfrac13\pi R^2 H = \tfrac13\pi(\qty{15e3}{m})^2(\qty{3e3}{m}) =
\qty{7.07e11}{m^3}$, y $\Delta m = (2800 - 1030)V = \qty{1.25e15}{kg}$.
(b) El centro de masa está a \qty{0.75}{km} sobre el fondo, a
$d = 5 - 0.75 = \qty{4.25}{km}$ bajo la superficie del mar. Allí la anomalía de potencial es
$T = -G\Delta m/d$, así que $N = G\Delta m/(\gamma d) =
\num{6.674e-11}\times\num{1.25e15}/(9.81\times4250) = \qty{2.0}{m}$.
(c) La mayor parte de la masa del cono está en su parte inferior, ancha, a entre 5 y
\qty{16}{km} del punto sobre la cima: el promedio de $1/r$ sobre la masa es menor que
$1/\qty{4.25}{km}$, y la masa puntual sobrestima el potencial. Sumando discos delgados resulta
unos \qty{1.1}{m}.
(d) Una raíz de \omterm{def:b1:layered-earth:layers}{corteza} ligera debajo pondría un déficit de la misma masa más profundo;
los dos potenciales casi se cancelan, y el abultamiento sería mucho menor, pero no nulo,
ya que el déficit está más lejos de la superficie que el exceso.
\end{solution}

\begin{solution}{exo:b1:isostasy-gravity:12}
(a) La velocidad de levantamiento es $-\mathrm dw/\mathrm dt = w_0\,e^{-t/\tau}/\tau = w/\tau$.
(b) $w = 818\,e^{-10000/4000} = 818\times0.0821 = \qty{67.1}{m}$, y la velocidad es
$\qty{67.1}{m}/\num{4000}\,\unit{yr} = \qty{16.8}{mm.yr^{-1}}$, más de una vez y
media el valor medido. (c) Con $\tau = 3100$ años, $w =
818\,e^{-3.226} = \qty{32.5}{m}$ y la velocidad es \qty{10.5}{mm.yr^{-1}}; con
$\tau = \num{68000}$ años, $w = 818\,e^{-0.147} = \qty{706}{m}$ y la velocidad
es \qty{10.4}{mm.yr^{-1}}. Los dos implican hundimientos restantes muy distintos.
Un hundimiento no compensado es un déficit de masa de roca del \omterm{def:b1:layered-earth:layers}{manto}: sobre una región amplia
da una \omterm{def:b1:isostasy-gravity:anomalies}{anomalía de aire libre} de unos $-2\pi G\rho_m w$, es decir
$-0.0419\times3.3\times706 \approx \qty{-98}{\milli\gal}$ para el caso lento,
frente a \qty{-4}{\milli\gal} para el rápido. Un levantamiento gravimétrico (o la altura de
antiguas líneas de costa) los distingue. (d) La constante de tiempo mide con qué rapidez
el \omterm{def:b1:layered-earth:layers}{manto} vuelve a fluir bajo el hundimiento: su viscosidad.
\end{solution}

\begin{solution}{pb:b1:isostasy-gravity:1}
\textbf{1.} $r = h\rho_c/(\rho_m - \rho_c) = 4\times2700/600 = \qty{18}{km}$.

\textbf{2.} La \omterm{def:b1:layered-earth:layers}{corteza} tiene $35 + 4 + 18 = \qty{57}{km}$ de espesor; su base, el \omterm{def:b1:layered-earth:discontinuity}{Moho},
está a $35 + 18 = \qty{53}{km}$ bajo el nivel del mar.

\textbf{3.} Meseta: \qty{57}{km} de \omterm{def:b1:layered-earth:layers}{corteza} y $60 - 53 = \qty{7}{km}$ de \omterm{def:b1:layered-earth:layers}{manto},
$\num{57000}\times2700 + \num{7000}\times3300 = \num{1.539e8} + \num{2.31e7} =
\qty{1.770e8}{kg.m^{-2}}$. Llanura: $\num{35000}\times2700 + \num{25000}\times3300
= \num{9.45e7} + \num{8.25e7} = \qty{1.770e8}{kg.m^{-2}}$. Son iguales, como
\cref{prop:b1:isostasy-gravity:equal-mass} exige.

\textbf{4.} $p = g\times\text{masa por unidad de área} = 9.81\times\num{1.770e8} =
\qty{1.74e9}{Pa}$, unos \qty{1.7}{GPa}, bajo ambas columnas.

\textbf{5.} La raíz debe compensar, con su menor densidad, la masa de la meseta: cada metro de raíz quita $\rho_m - \rho_c = \qty{600}{kg.m^{-3}}$
de masa, y cada metro de meseta añade \qty{2700}{kg.m^{-3}}, así que $2700/600 = 4.5$
metros de raíz son necesarios. Con una \omterm{def:b1:layered-earth:layers}{corteza} de \qty{2900}{kg.m^{-3}}, la razón
sería $2900/400 = 7.25$: cuanto menor es el contraste, más profunda es la raíz.

\textbf{6.} El \omterm{def:b1:layered-earth:discontinuity}{Moho}, el fuerte aumento de la velocidad sísmica, se encontraría más profundo bajo el borde de la meseta, desde
\qty{35}{km} bajo la llanura hasta unos \qty{53}{km} bajo la meseta.

\textbf{7.} $\rho = 2700\times100/104 = \qty{2596}{kg.m^{-3}}$.

\textbf{8.} $\num{104000}\times2596.2 = \qty{2.700e8}{kg.m^{-2}}$ y
$\num{100000}\times2700 = \qty{2.700e8}{kg.m^{-2}}$.

\textbf{9.} $2700 - 2596 = \qty{104}{kg.m^{-3}}$, o $104/2700 =
\qty{3.85}{\percent}$.

\textbf{10.} $\Delta\rho/\rho = \alpha\,\Delta T$ da $\Delta T =
0.0385/\num{3e-5} = \qty{1280}{K}$ como promedio sobre toda la columna de \qty{104}{km}. Eso fundiría buena parte de la columna: el calor
por sí solo no puede hacer esta meseta lo bastante ligera, y el mecanismo de Pratt no puede ser aquí el principal.

\textbf{11.} Las columnas de Pratt no tienen raíz: el \omterm{def:b1:layered-earth:discontinuity}{Moho} se queda a unos \qty{35}{km} bajo la meseta, como bajo la llanura.

\textbf{12.} \qty{52}{km} está cerca de los \qty{53}{km} de Airy y lejos de los \qty{35}{km} de Pratt: la meseta está compensada sobre todo por una raíz cortical.

\textbf{13.} $0.0419\times2.7\times4000 = \qty{452.5}{\milli\gal}$.

\textbf{14.} Sin compensar (\cref{prop:b1:isostasy-gravity:broad}):
$\Delta g_{\mathrm{FA}} \approx \qty{+452.5}{\milli\gal}$,
$\Delta g_{\mathrm B} \approx 0$.

\textbf{15.} Compensada: $\Delta g_{\mathrm{FA}} \approx 0$ y
$\Delta g_{\mathrm B} \approx \qty{-452.5}{\milli\gal}$. La masa que falta en la raíz
da $-0.0419\times0.6\times\num{18000} = \qty{-452.5}{\milli\gal}$, la misma,
ya que $\Delta\rho\,r = \rho_c h$.

\textbf{16.} $\sin^2\ang{30} = 0.25$ y $\sin^2\ang{60} = 0.75$:
$\gamma = 9.780327\times(1 + 0.0013256 - 0.0000044) = \qty{9.793249}{m.s^{-2}}$.
Entonces $g - \gamma = 9.781085 - 9.793249 = \qty{-1216.4}{\milli\gal}$; la \omterm{def:b1:isostasy-gravity:anomalies}{corrección de aire libre}
es $0.3086\times4000 = \qty{1234.4}{\milli\gal}$, así que
$\Delta g_{\mathrm{FA}} = \qty{+18.0}{\milli\gal}$, y
$\Delta g_{\mathrm B} = 18.0 - 452.5 = \qty{-434.5}{\milli\gal}$.

\textbf{17.} $434.5/452.5 = 0.96$: la meseta está compensada en unos \qty{96}{\percent}
(la pequeña \omterm{def:b1:isostasy-gravity:anomalies}{anomalía de aire libre} positiva es el \qty{4}{\percent} sin compensar).

\textbf{18.} No. Cerca del centro de una meseta amplia, cada capa anómala actúa
como una placa, cuya atracción no depende de su profundidad: solo importa el déficit de masa
por unidad de área, que es el mismo, $\rho_c h$, en ambos modelos. Los dos difieren
solo cerca de los bordes, donde la profundidad de las masas compensadoras cambia
la forma de la anomalía.

\textbf{19.} Lejos del borde, a cada lado, la \omterm{def:b1:isostasy-gravity:anomalies}{anomalía de aire libre} es casi cero.
A unos pocos kilómetros dentro del borde, la meseta bajo la estación está cerca y se extiende lejos solo a un lado: sigue tirando con casi todo su valor de placa.
La raíz, \qty{35}{km} o más abajo, se ve desde allí como un cuerpo que termina casi recto por debajo: tira con algo más de la mitad de su efecto. La diferencia es un pico positivo. A unos pocos kilómetros fuera, la meseta tira muy poco, mientras que la raíz sigue tirando con casi la mitad de su efecto: un valle negativo. El perfil es un par positivo--negativo a cada lado del borde, como en la figura del capítulo.

\textbf{20.} La \omterm{def:b1:layered-earth:layers}{corteza} es ahora de $57 - 1 = \qty{56}{km}$ de espesor, $\qty{21}{km}$
más que la de la llanura. En equilibrio este exceso se reparte como $h' + r'$ con
$r' = 4.5\,h'$: $h' = 21/5.5 = \qty{3.818}{km}$ y $r' = \qty{17.18}{km}$.

\textbf{21.} La meseta se alza ahora \qty{3818}{m}.

\textbf{22.} La roca que ahora está en la superficie estaba \qty{1}{km} bajo la antigua superficie,
a \qty{3000}{m}; ahora está a \qty{3818}{m}: ha subido \qty{818}{m}, es decir $1000\times\rho_c/\rho_m$.

\textbf{23.} Cada kilómetro erosionado baja la superficie $1 - 2700/3300 =
0.182$ de kilómetro, así que bajarla \qty{4}{km} exige
$4/0.182 = \qty{22}{km}$ de erosión, todo el exceso de espesor $h + r$. En un
cinturón de montañas antiguo y desgastado, las rocas en la superficie estuvieron, por tanto, una vez a unos
veinte kilómetros de profundidad: las raíces profundas del cinturón están expuestas.

\textbf{24.} La superficie baja $1000 - 818$: por cada kilómetro de roca erosionada, el descenso neto de la superficie es
\textbf{\qty[text-series-to-math=true]{182}{m}}.
\end{solution}
